%%%%%%%%%%%%%%%%%%%%%%%%%%%%%%%%%%%%%%%%%%%%%%%%%%%%%%%%%%%%%%%%%%%%%%%%%%%%%%%%
%   Copyright 2009, Oracle and/or its affiliates.
%   All rights reserved.
%
%
%   Use is subject to license terms.
%
%   This distribution may include materials developed by third parties.
%
%%%%%%%%%%%%%%%%%%%%%%%%%%%%%%%%%%%%%%%%%%%%%%%%%%%%%%%%%%%%%%%%%%%%%%%%%%%%%%%%

\chapter{Overloading and Multiple Dispatch}
\chaplabel{multiple-dispatch}

\note{Keyword and varargs parameters are not yet supported.}

Fortress allows functions and methods (collectively called \emph{functionals})
to be \emph{overloaded}.  That is,
there may be multiple declarations for the same functional name
visible in a single scope (which may include inherited method declarations),
and several of them may be applicable to any particular functional call.

However, with these benefits comes the potential for ambiguous calls at run time.
For an overloaded functional call,
the most specific applicable declaration is chosen, if any.
Otherwise, any of the applicable declarations such that
no other applicable declaration is more specific than them is chosen.


In this chapter, we describe how to determine which declarations
are \emph{applicable} to a particular functional call,
and when several are applicable, how to select among them.
We also provide overloading rules for the declarations of functionals
to eliminate \emph{some possibility} of ambiguous calls at run time,
whether or not these calls actually appear in the program.
\secref{overloading-terms} introduces some terminology and notation.
In \secref{functional-applicability},
we show how to determine which declarations
are applicable to a \emph{named functional call}
(a function call described in \secref{function-calls}
or a naked method invocation described in \secref{naked-method-calls})
when all declarations have only ordinary parameters
(without varargs or keyword
parameters).
We discuss how to handle dotted method calls
(described in \secref{dotted-method-calls})
in \secref{dotted-method-applicability},
and declarations with varargs and keyword
parameters in \secref{applicability-with-special-params}.
Determining which declaration is applied,
if several are applicable,
is discussed in \secref{resolving-overloading}.


\section{Principles of Overloading}
\seclabel{overloading-terms}

Fortress allows multiple functional declarations of the same name to
be declared in a single scope.  However, recall from \chapref{names}
the following shadowing rules:
\begin{itemize}
\item
dotted method declarations shadow top-level function declarations with
the same name, and
\item
dotted method declarations provided by a trait or object declaration
or object expression
shadow functional method declarations with the same name that are
provided by a different trait or object declaration
or object expression.
\end{itemize}
Also, note that a trait or object declaration
or object expression must not have a
functional method declaration and a dotted method declaration with the
same name, either directly or by inheritance.
Therefore, top-level functions can overload
with other top-level functions and functional methods,
dotted methods with other dotted methods,
and functional methods with other functional methods and top-level functions.
It is a static error if any top-level function declaration is
more specific than any functional method declaration.
If a top-level function declaration is overloaded with a functional method
declaration, the top-level function declaration must not be more specific than
the functional method declaration.


Operator declarations with the same name but different fixity
are not a valid overloading; they are unambiguous declarations.
An operator method declaration whose name is one of the operator parameters
(described in \secref{opr-ident})
of its enclosing trait or object may be overloaded with other operator
declarations in the same component.
Therefore, such an operator method declaration
must satisfy the overloading rules
(described in \chapref{overloaded-declarations})
with every operator declaration in the same component.

\note{This restriction will be relaxed.}

Declarations with the same functional name may have different static
parameters, or some of them may have static parameters and others none.
A declaration with static parameters is read as one declaration whose
parameter type ranges over every instance of it whose static arguments
satisfy their bounds: it is applicable to a call when some such instance is
(\secref{functional-applicability}), and it is compared with other
declarations on the parameter types they declare, taken over all their
instances (\secref{resolving-overloading}).
A declaration without static parameters is the case with none.
The rules that a set of such declarations must satisfy are given in
\chapref{overloaded-declarations}.
\revision{revival-overloading}{The Working Draft of February 2011 said here that
declarations of one name may not differ in their static parameters, nor may one
have static parameters and another not, so that static parameters do not enter
into applicability.
Neither implementation enforced that: the interpreter has overloaded
declarations with different static parameters since 2007, and the compiled type
checker applies to them the rules of the type group's paper of 2011, which reads
the parameter type of a declaration with static parameters as quantified over
them; \library\ overload across different static parameters throughout.
This chapter and \chapref{overloaded-declarations} now state that model.}



Recall from \chapref{types} that we write
$T \Ovrsubtype U$ when $T$ is a subtype of $U$,
and $T \strictsubtype U$ when $T \Ovrsubtype U$ and $T \neq U$.


\section{Applicability to Named Functional Calls}
\seclabel{functional-applicability}
In this section, we show how to determine which declarations are
applicable to a named functional call when all declarations have only
ordinary parameters (i.e., neither varargs nor keyword parameters).

For the purpose of defining applicability, a named functional call can
be characterized by the name of the functional and its argument type.
Recall that a functional has a single parameter, which may be a tuple
(a dotted method has a receiver as well).
 We abuse notation by using
\emph{static call $\f(\As)$} to refer to a named functional call with
name $\f$ and whose argument has static type $\As$, and
\emph{dynamic call $\f(\Xs)$} to refer to a named functional call with
We use \emph{call $\f(\Cs)$} to refer to a named functional call with
name $\f$ and whose argument, when evaluated, has dynamic type $\Cs$.
(Note that if the type system is sound---and we certainly hope that
it is!---then $\Xs \Ovrsubtype \As$ for all well-typed calls to $\f$.)
We use the term \emph{call $\f(\Cs)$} to refer to static and dynamic calls
collectively.
We assume throughout this chapter that the static arguments that a call writes
(\secref{idn-ref}, \secref{dotted-method-calls}) have been substituted for the
static parameters of the declarations that they instantiate.
The static arguments that a call does not write are inferred for the
declaration that the call selects, once it is selected
(\chapref{type-inference}); until then a declaration with static parameters is
taken over all its instances.
\revision{revival-overloading}{The Working Draft of February 2011 assumed here
that all static variables in functional calls have been instantiated or
inferred, which with the sentences below had the static arguments a call does not
write inferred before the declarations are compared.}


We also use \emph{function declaration $\f(\Ps):\Us$} to refer to a
function declaration with function name $\f$, parameter type $\Ps$,
and return type $\Us$.


For method declarations, we must take into account the self parameter,
as follows:

A \emph{dotted method declaration $\Ps_0.\f(\Ps):\Us$} is a dotted
method declaration with name $f$, where $\Ps_0$ is the trait or object
type in which the declaration appears, $\Ps$ is the parameter type,
and $\Us$ is the return type.  (Note that despite the suggestive
notation, a dotted method declaration does not explicitly list its
self parameter.)

A \emph{functional method declaration $\f(\Ps):\Us$ with self
parameter at $i$} is a functional method declaration with name $f$,
with the parameter \EXP{\mathord{\KWD{self}}} in the $i$th position
of the parameter type $\Ps$, and return type $\Us$.  Note that the static
type of the self parameter is the trait or object trait type in which
the declaration $\f(\Ps):\Us$ occurs.  In the following, we will use
$\Ps_i$ to refer to the $i$th element of $\Ps$.

We elide the return type of a declaration, writing $\f(\Ps)$ and
$\Ps_0.\f(\Ps)$, when the return type is not relevant to the
discussion.
Note that static parameters may appear in the types $\Ps_0$, $\Ps$, and $\Us$.


A declaration $\f(\Ps)$ is \emph{applicable} to a call $\f(\Cs)$
if the call is in the scope of the declaration and $\Cs \Ovrsubtype \Ps$.
(See \chapref{names} for the definition of scope.)
If the declaration has static parameters, it is applicable to the call if some
instance of it whose static arguments satisfy their bounds is, that is, if
$\Cs \Ovrsubtype \Ps'$ for the parameter type $\Ps'$ of such an instance.
Its static arguments are not inferred to decide this: they are inferred only for
the declaration that the call selects, after it is selected
(\secref{static-arg-inference}).
\revision{revival-overloading}{The Working Draft of February 2011 said that
static parameters in the parameter type are inferred as described in
\chapref{type-inference} before the applicability of the declaration to the call
is checked.
The type group's paper of 2011 defines a declaration with static parameters to
be applicable when at least one of its instances is.}



Note that a named functional call $\f(\Cs)$ may invoke a dotted method
declaration if the declaration is provided by the trait or object
enclosing the call.  To account for this, let $\Cs_0$ be the trait or
object declaration immediately enclosing the call.
%
Then we consider a named functional call $\f(\Cs)$ as $\Cs_0.\f(\Cs)$
if $\Cs_0$ provides dotted method declarations applicable to $\f(\Cs)$, and
use the rule for applicability to dotted method calls (described in
\secref{dotted-method-applicability}) to determine which declarations
are applicable to $\Cs_0.\f(\Cs)$.



\section{Applicability to Dotted Method Calls}
\seclabel{dotted-method-applicability}

Dotted method applications can be characterized similarly to named
functional applications, except that, analogously to dotted method declarations,
we use $\Cs_0$ to denote the dynamic type of the receiver object,
and, as for named functional calls, $\Cs$ to denote the dynamic
type of the argument of a  dotted method call.
We write $\Cs_0.\f(\Cs)$ to refer to the  call.

A dotted method declaration $\Ps_0.\f(\Ps)$ is \emph{applicable}
to a dotted method call $\Cs_0.\f(\Cs)$ if $\Cs_0 \Ovrsubtype
\Ps_0$ and $\Cs \Ovrsubtype \Ps$.
If the declaration has static parameters, it is applicable to the call if some
instance of it whose static arguments satisfy their bounds is, that is, if
$\Cs_0 \Ovrsubtype \Ps_0'$ and $\Cs \Ovrsubtype \Ps'$ for the types $\Ps_0'$
and $\Ps'$ of such an instance; its static arguments are inferred only for the
declaration that the call selects (\secref{static-arg-inference}).
\revision{revival-overloading}{The Working Draft of February 2011 said, as for
named functional calls, that static parameters in $\Ps_0$ and $\Ps$ are inferred
before the applicability of the declaration is checked.}



\section{Applicability for Functionals with Varargs and Keyword Parameters}
\seclabel{applicability-with-special-params}

The basic idea for handling varargs and keyword
parameters is that we
can think of a functional declaration that has such parameters as
though it were (possibly infinitely) many declarations, one for each
set of arguments it may be called with.  In other words, we expand
these declarations so that there exists a declaration for each number
of arguments that can be passed to it.


A declaration with a varargs parameter corresponds to an infinite
number of declarations, one for every number of arguments that may be
passed to the varargs parameter.  In practice, we can bound that
number by the maximum number of arguments that the functional is
called with anywhere in the program (in other words, a given program
will contain only a finite number of calls with different numbers of
arguments).  The expansion described here is a conceptual one to
simplify the description of the semantics; we do not expect a real
implementation to actually expand these declarations at compile time.
%
For example, the following declaration:
%f(x : ZZ, y : ZZ, z : ZZ...):ZZ
\begin{Fortress}
\(f(x \mathrel{\mathtt{:}} \mathbb{Z}, y \mathrel{\mathtt{:}} \mathbb{Z}, z \mathrel{\mathtt{:}} \mathbb{Z}\ldots)\COLONOP\mathbb{Z}\)
\end{Fortress}
would be expanded into:
%f(x : ZZ, y : ZZ):ZZ
%f(x : ZZ, y : ZZ, z1 : ZZ):ZZ
%f(x : ZZ, y : ZZ, z1 : ZZ, z2 : ZZ):ZZ
%f(x : ZZ, y : ZZ, z1 : ZZ, z2 : ZZ, z3 : ZZ):ZZ
%...
\begin{Fortress}
\(f(x \mathrel{\mathtt{:}} \mathbb{Z}, y \mathrel{\mathtt{:}} \mathbb{Z})\COLONOP\mathbb{Z}\)\\
\(f(x \mathrel{\mathtt{:}} \mathbb{Z}, y \mathrel{\mathtt{:}} \mathbb{Z}, z_{1} \mathrel{\mathtt{:}} \mathbb{Z})\COLONOP\mathbb{Z}\)\\
\(f(x \mathrel{\mathtt{:}} \mathbb{Z}, y \mathrel{\mathtt{:}} \mathbb{Z}, z_{1} \mathrel{\mathtt{:}} \mathbb{Z}, z_{2} \mathrel{\mathtt{:}} \mathbb{Z})\COLONOP\mathbb{Z}\)\\
\(f(x \mathrel{\mathtt{:}} \mathbb{Z}, y \mathrel{\mathtt{:}} \mathbb{Z}, z_{1} \mathrel{\mathtt{:}} \mathbb{Z}, z_{2} \mathrel{\mathtt{:}} \mathbb{Z}, z_{3} \mathrel{\mathtt{:}} \mathbb{Z})\COLONOP\mathbb{Z}\)\\
\(\ldots\)
\end{Fortress}
A declaration with a varargs parameter is applicable to a call if
any one of the expanded declarations is applicable.



\section{Overloading Resolution}
\seclabel{resolving-overloading}

\note{
Victor: Does this paragraph, other than the last sentence,
really belong in this section?}

To evaluate a given functional call,
it is necessary to determine which functional declaration
to dispatch to.
To do so,
we consider the declarations that are applicable to that call at run time.
If there is exactly one such declaration,
then the call dispatches to that declaration.
If there is no such declaration,
then the call is \emph{undefined}, which is a static error.
(However, see \secref{applicability-with-coercion}
for how coercion may add to the set of applicable declarations.)
If multiple declarations are applicable to the call at run time,
then we choose an arbitrary declaration among the declarations such that
no other applicable declaration is more specific than them.
A declaration with static parameters is applicable to a call at run time, as
to a static call, when some instance of it whose static arguments satisfy
their bounds is applicable to the call, so the declaration to which a call
dispatches may be one that the static call did not select.
A declaration that the static call selected runs at the instance inferred for
the call (\secref{static-arg-inference}); any other declaration runs at the
instance in which each type parameter is the intersection of the upper bounds
placed on it: its declared bound, \TYP{Any} if it has none
(\secref{type-param}), and the bounds that the dynamic types of the arguments
place on it through the parameter types and that the static type of the call
places on it through the return type.
Where the type parameter is a static argument of a parameterized trait in a
parameter type, as \VAR{X} is in \EXP{t\COLON\TYP{Tag}\llbracket{}X\rrbracket},
the value of the argument fixes it, since no value belongs to two
instantiations of one parameterized trait (\secref{trait-types}); where it is
the whole type of a parameter, the dynamic type of the argument is only a
lower bound, and the type parameter is its bound under the static type of the
call.
The type parameters of a declaration that only dispatch reaches are never
instantiated with \TYP{BottomType}, nor with the union of the types of the
arguments.
\revision{revival-dispatch}{The Working Draft of February 2011 did not say at
which instance a declaration with static parameters runs when dispatch reaches
it and the static call selected another.
The rule stated is the one by which the dispatcher of the type group's paper of
2019 instantiates a declaration (Park, Hong, Steele and Ryu, ``Polymorphic
Symmetric Multiple Dispatch with Variance'', POPL 2019); the team had assumed
the union of the lower bounds and could not prove it correct.
Both implementations of the revival take the dynamic type of the argument where
the value does not fix the type parameter: the interpreter because it has no
static types, and so no static type of the call to bound the instance, as with
its choice of coercions (\secref{resolving-coercion}); the compiled path because
its dispatcher reads the static arguments of a dispatched declaration from the
value and is not given the static type of the call.
Where the value fixes the type parameter, the two rules give the same instance.
The compiled path's departure is gated as an expected failure (row~538 of the
revival's gap ledger), and the compiled path does not yet run the second
example below (row~496).}



We use the subtype relation to compare parameter types to determine
a more specific declaration.
Formally, a declaration $\f(\Ps)$ is \emph{more specific} than a declaration
$\f(\Qs)$ if $\Ps \strictsubtype \Qs$.  Similarly, a declaration
$\Ps_0.\f(\Ps)$ is more specific than a declaration $\Qs_0.\f(\Qs)$
if $\Ps_0 \strictsubtype \Qs_0$ and $\Ps \strictsubtype \Qs$.
(See \secref{resolving-coercion} for how coercion
changes the definition of ``more specific''.)
Restrictions on the definition of overloaded functionals
(see \chapref{overloaded-declarations})
guarantee that among all applicable declarations,
one is more specific than all the others.
Where declarations have static parameters, they are compared on the parameter
types they declare, each taken over all its instances: a declaration is more
specific than another when the second is applicable to every call to which the
first is applicable, and not the other way round.
Where neither has static parameters, this is the subtype relation above.
No static arguments are inferred for the comparison.
\revision{revival-overloading}{The Working Draft of February 2011 said that
static parameters are inferred as described in \chapref{type-inference} before
the parameter types are compared.
The definition stated is that of the type group's paper of 2011, which checks
it as subtyping of the parameter types read as existential types.}

For example, where \TYP{Shape} comprises the two object types \TYP{Circle}
and \TYP{Square}, of the declarations
%describe[\T extends Shape\](x: T): String = "generic"
%describe(x: Circle): String = "circle"
\begin{Fortress}
\(\VAR{describe}\llbracket{}T \KWD{extends} \TYP{Shape}\rrbracket(x\COLON{}T)\COLON \TYP{String} = \hbox{\rm``\STR{generic}''}\)\\
\(\VAR{describe}(x\COLON \TYP{Circle})\COLON \TYP{String} = \hbox{\rm``\STR{circle}''}\)
\end{Fortress}
the second is more specific than the first, although the first has an instance
whose parameter type is \TYP{Circle}.
The calls \EXP{\VAR{describe}(\TYP{Circle})} and \EXP{\VAR{describe}(s)}, where
\VAR{s} has static type \TYP{Shape} and holds \TYP{Circle}, run the second
declaration, and \EXP{\VAR{describe}(\TYP{Square})} runs the first.
And of the declarations
%grab[\X\](t: Tag[\X\]): Marker[\X\] = Marker[\X\]
%grab(a: Any): Any = Red
\begin{Fortress}
\(\VAR{grab}\llbracket{}X\rrbracket(t\COLON \TYP{Tag}\llbracket{}X\rrbracket)\COLON \TYP{Marker}\llbracket{}X\rrbracket = \TYP{Marker}\llbracket{}X\rrbracket\)\\
\(\VAR{grab}(a\COLON \TYP{Any})\COLON \TYP{Any} = \TYP{Red}\)
\end{Fortress}
where \TYP{Marker} is an object with a type parameter and \TYP{Red} an object
type, the call \EXP{\VAR{grab}(a)}, where \VAR{a} has static type \TYP{Any},
selects the second declaration statically; if \VAR{a} holds a value of an
object type that extends \EXP{\TYP{Tag}\llbracket\TYP{Red}\rrbracket}, the
call dispatches to the first, which is applicable to that value and more
specific, at \EXP{X = \TYP{Red}}, which the value fixes, and returns
\EXP{\TYP{Marker}\llbracket\TYP{Red}\rrbracket}.
\revision{revival-dispatch}{The Working Draft of February 2011 forbade both
pairs of declarations, which differ in their static parameters.}
